\section{Related Work}
\label{sec:related}

Our study connects three bodies of work that have largely developed
separately: driving-gaze datasets, validation studies of wearable
accuracy, and methods that correct or qualify gaze. We review each in
turn and state what it establishes for gaze supervision in real driving.

\textbf{Gaze supervision and driving datasets.}
Human gaze is an attractive supervisory signal for driving. It has been
used as an input or auxiliary target for driving models~\cite{liu2019gazeModel,makrigiorgos2019attentionDriving,abdelkarim2021ggnet,xia2020periphery}
and to reduce causal confusion in learned agents~\cite{biswas2024gazeSupervision}.
The underlying recordings come from three kinds of sources: wearable
trackers in real vehicles, as in DR(eye)VE and Look Both
Ways~\cite{palazzi2019dreyeve,kasahara2022lookBothWays}; observers
watching recorded driving video in the laboratory, as in BDD-A, MAAD and
DADA~\cite{xia2018bdda,gopinath2021maad,fang2022dada}; and driving
simulators~\cite{contreras2026laia,govil2026drivergaze360}. Together
they have driven substantial progress in attention modelling.

The field is also aware that these labels are imperfect. Gaze has been
described as a noisy approximation of the causally relevant image
regions~\cite{biswas2024gazeSupervision}, and existing gaze maps as noisy
and partly irrelevant~\cite{liang2023cueing}. Systematic audits make the
sources of this noise concrete~\cite{kotseruba2024dataLimitations,kotseruba2024taskContext}.
In DR(eye)VE, videos lacked timestamps and were typically misaligned by
several frames; homography projection of wearable gaze into the roof
camera produced outliers in more than 13\% of frames; and blinks and
saccades make up over a quarter of the samples. The raw tracker output of
Look Both Ways is not available for independent re-analysis.
What these audits could not assess is the angular error of the gaze
itself. In the datasets we examined, the tracker output, once mapped into
the scene, \emph{is} the label, and its error against independently
identified physical targets during driving is not reported. DR(eye)VE
measured tracker error on calibration patterns at 5--15\,m to set the
width of its fixation maps, but the resulting value is not given in the
dataset paper~\cite{palazzi2019dreyeve}; simulator datasets assume a small
device error~\cite{contreras2026laia}; and on-road accuracy figures that do
exist are calibration-time values for remote trackers~\cite{stapel2022driverPerception}.
A downstream gain from gaze supervision therefore does not by itself
show that the labels are accurate.

This motivates our first contribution. Drivers announce fixations on
traffic objects; the targets are localised and manually verified in the
eye-tracker scene view, and each reference is synchronised with vehicle
motion, LiDAR-guided range and scenario labels. Where earlier datasets
provide gaze, our benchmark provides gaze together with a reference
against which it can be measured.

\textbf{Wearable eye-tracking accuracy and robustness.}
Measuring wearable accuracy against known targets is well established,
with validated protocols and tools~\cite{niehorster2024glassesvalidator,niehorster2026fundamentals7}
and studies of viewing distance and gaze angle~\cite{macinnes2018wearable},
body motion~\cite{hooge2023robust} and walking~\cite{onkhar2024evaluating}
on earlier devices. For Pupil Labs Neon, whose calibration-free gaze
estimation makes natural driving recordings practical, the evidence is
growing quickly. The manufacturer's report covers 206 laboratory and 402
everyday-environment participants, with median accuracy of about
1.8\textdegree{} without and 1.3\textdegree{} with a person-specific
offset at 1.3\,m~\cite{baumann2026neonAccuracy}. Independent studies have
probed individual influences: head pose, temporal stability and external
infrared sources~\cite{foucher2026neon}, gaze position and background
luminance~\cite{alecka2026neon}, pupil size~\cite{salari2026pupilSize,maquiling2026vergencePupil}
and glasses slippage~\cite{pettenella2026slippage}.

Taken together, these studies paint an encouraging picture: offset-corrected
accuracy of roughly 1.3--1.5\textdegree{} in the laboratory and mostly
modest condition effects. They also reveal specific sensitivities, and
not all findings agree. A black background reduced accuracy in one
study~\cite{alecka2026neon}, whereas another found no accuracy change between bright and dark
screens but did report an apparent gaze shift of almost one
degree~\cite{salari2026pupilSize}. The manufacturer reports lower accuracy in the lower visual field and
outdoors~\cite{baumann2026neonAccuracy}, and infrared light from a second tracker degraded Neon in a
between-group comparison~\cite{foucher2026neon}.

Nearly all of these evaluations, however, use screens, posters or
predefined targets at no more than a few metres, with the participant
seated, standing or walking. Extrapolation to larger depths is motivated
by the small change of vergence beyond 1.3\,m~\cite{baumann2026neonAccuracy}, which is an argument
rather than a measurement. In real vehicles, Neon has so far been
reported for signal availability in flight~\cite{xu2026signalConstruct},
and wearable tracking in urban driving for behavioural measures rather
than accuracy~\cite{calvi2023comparing}. Yet in driving the relevant objects are tens of metres away, where even a
small angular error becomes a substantial lateral displacement.
We therefore complement existing validations in two ways: a real-driving
benchmark that measures accuracy where gaze supervision is actually
collected, and a controlled repeated-measures study whose factors
(viewing distance, position in the field of view, pupil size, indoor and
outdoor illumination, static and pursuit fixation) are chosen because they
occur in driving. Rather than questioning the device, our experiments ask
under which conditions its output can serve as supervision directly and
under which it should be verified.

\textbf{In-situ correction and quality assessment.}
A recurring observation in these studies is that a large share of the
error is a wearer-specific offset. Correcting it is established practice:
Neon provides a one-point person-specific offset whose benefit has been
reported at scale~\cite{baumann2026neonAccuracy} and
independently~\cite{foucher2026neon}, and several methods infer
corrections from natural viewing instead of dedicated targets.
Implicit calibration derives probable fixation targets from saliency or
expert labels and stabilises after a few dozen fixations on
screens~\cite{kasprowski2019implicitCalibration,chang2019salgaze}, and a
handful of in-vehicle samples personalises remote driver-facing gaze
models~\cite{nagpure2025learningToAdapt}.
Closest to our work, Jia et al.~\cite{jia2018gaze} corrected the offset of
wearable glasses using detected stop signs. In a moving car, three
passengers were instructed to follow the signs, targets were assigned by gaze proximity, and errors
were reported in pixels. This establishes the feasibility of
road-object-based correction. In the studies we reviewed, however, two questions that matter for
supervision data remain open. First, does a correction transfer
across conditions? A mismatch between calibration and use distance
measurably raises error for head-mounted trackers~\cite{barz2016prediction},
and Neon offsets in laboratory studies are set on near-field screen targets
under a single lighting condition~\cite{baumann2026neonAccuracy,alecka2026neon}. We found no comparison
of laboratory-derived and in-situ offsets on identical held-out driving
targets, nor an analysis of how many references are needed and how they
should be arranged. Second, does added modelling help beyond a simple offset?
Simple regressions are often hard to beat for gaze-error
modelling~\cite{kar2020mlgaze}, while pupil size~\cite{choe2016pupilSize},
auxiliary motion signals~\cite{park2026motionCorrected} and slippage
robustness~\cite{santini2019grip} suggest where targeted extensions could help.

Beyond correction, supervision data need to be annotated for the
tasks they will serve. Building blocks exist: quality gates on fixation
stability, head motion and blinks~\cite{baumann2026neonAccuracy}, eyelid
and eye-opening estimation~\cite{fuhl2024pistol,zhang2026pupeyes},
image-quality factors that rank angular error for appearance-based
estimators~\cite{jamalifard2026fiud}, and a recalibration threshold in
an on-road study~\cite{stapel2022driverPerception}. Whether data are
good enough depends on the research question~\cite{niehorster2026fundamentals7};
in-cabin fixations and blinks, for example, are noise for scene-saliency
learning but not invalid gaze in general~\cite{kotseruba2024dataLimitations}.

We build on these precedents rather than claim in-situ correction as a new
concept. Our contribution is a systematic transfer comparison: laboratory-derived
and driving-derived offsets evaluated on the same held-out driving targets
in angular units, the number and arrangement of references required, and
targeted individual spatial and temporal extensions evaluated against a
strong median-offset baseline. We also connect correction to acquisition.
The system that generates verified references annotates samples with
task-aware suitability flags that combine viewing region, scene sharpness
and motion, eye visibility and gaze dynamics. Unlike error predictors, these
flags are transparent annotations for selection by downstream tasks, not
estimates of instantaneous accuracy, and sparse verified fixations provide
an in-situ offset check.

\textbf{Summary.}
The reviewed work shows why gaze supervision is valuable, how accurate
wearable trackers are under controlled conditions, and which corrections
and quality indicators are available. It does not connect these threads in
real driving. We link them: verified in-situ references make gaze accuracy
measurable in actual driving; a driving-motivated controlled study
characterises where Neon is robust; and sparse condition-matched
correction with quality-annotated acquisition turns these findings into
practice. The goal is gaze-supervision data whose measurement quality is
evaluated rather than assumed.